% Part: sets-functions-relations
% Chapter: relations
% Section: relations-as-sets

\documentclass[../../../include/open-logic-section]{subfiles}

\begin{document}

\olfileid[pt]{sfr}{rel}{set}
\olsection{Relações como Conjuntos}

\begin{explain}
Na \olref[sfr][set][imp]{sec}, mencionamos alguns conjuntos importantes:
$\Nat$, $\Int$, $\Rat$, $\Real$. Sem dúvida que se lembram de algumas
relações interessantes entre os !!{element}s de alguns destes conjuntos.
Por exemplo, cada um destes conjuntos tem uma \emph{relação de ordem}
completamente padrão sobre si. Há também a relação \emph{ser idêntico a},
segundo a qual cada objeto só se relaciona consigo próprio e com nenhum outro. Há
muitas outras relações interessantes que encontraremos, e ainda mais
relações possíveis. Antes de as revermos, contudo, começaremos por
observar que podemos ver relações como um tipo especial de conjunto. 
  
Para isso, recordemos duas coisas de \olref[sfr][set][pai]{sec}. Primeiro,
recordemos a noção de \emph{par ordenado}: dados $a$ e $b$, podemos
formar~$\tuple{a, b}$. É importante que a ordem dos elementos seja relevante
aqui. Assim, se $a \neq b$ então $\tuple{a, b} \neq \tuple{b, a}$.
(Contrastemos isto com pares não ordenados, O.s.\@, conjuntos de $2$ elementos, onde
$\{a, b\}=\{b, a\}$.) Segundo, recordemos a noção de \emph{produto
cartesiano}: se $A$ e $B$ são conjuntos, então podemos formar~$A \times B$, o
conjunto de todos os pares $\tuple{x, y}$ com $x \in A$ e $y \in B$. Em
particular, $A^{2}= A \times A$ é o conjunto de todos os pares ordenados
de~$A$.
  
Consideremos agora uma relação particular num conjunto: a relação~$<$
no conjunto~$\Nat$ dos números naturais. Seja o conjunto de todos os
pares de números $\tuple{n, m}$ onde $n<m$, O.s.\@,
\[
  R=\Setabs{\tuple{n, m}}{n, m \in \Nat \text{ e } n<m}.
\]
Há uma ligação estreita entre $n$ ser menor que $m$, e o
par $\tuple{n, m}$ ser membro de $R$, nomeadamente:
\[
      n<m\text{ se e só se }\tuple{n, m} \in R.
\]
De fato, sem qualquer perda de informação, podemos considerar o conjunto $R$
como \emph{sendo} a relação~$<$ em $\Nat$. 

Do mesmo modo podemos construir um subconjunto de $\Nat^{2}$ para qualquer
relação entre números. Reciprocamente, dado qualquer conjunto de pares de
números $S \subseteq \Nat^{2}$, há uma relação correspondente
entre números, nomeadamente, a relação segundo a qual $n$ se liga a $m$ se e só se
$\tuple{n, m} \in S$. Isto justifica a definição seguinte:
\end{explain}
  
\begin{defn}[Relação binária] 
Uma \emph{relação binária} num conjunto $A$ é um subconjunto de~$A^{2}$. Se $R
\subseteq A^{2}$ é uma relação binária em~$A$ e $x, y \in A$, por vezes
escrevemos $Rxy$ (ou $xRy$) para $\tuple{x, y} \in R$.
\end{defn}
  
\begin{ex}
  \ollabel{relations}
O conjunto $\Nat^{2}$ dos pares de números naturais pode ser listado numa
matriz bidimensional assim:
\[
  \begin{array}{ccccc}
  \mathbf{\tuple{ 0,0 }} & \tuple{ 0,1 } &
    \tuple{ 0,2 } & \tuple{ 0,3 } & \ldots\\
  \tuple{ 1,0 } & \mathbf{\tuple{ 1,1 }} &
    \tuple{ 1,2 } & \tuple{ 1,3 } & \ldots\\
  \tuple{ 2,0 } & \tuple{ 2,1 } &
    \mathbf{\tuple{ 2,2 }} & \tuple{ 2,3 } & \ldots\\
  \tuple{ 3,0 } & \tuple{ 3,1 } & \tuple{ 3,2 } &
    \mathbf{\tuple{ 3,3 }} & \ldots\\
  \vdots & \vdots & \vdots & \vdots & \mathbf{\ddots}
  \end{array}
\]
Colocamos a diagonal, aqui, em negrito, pois o subconjunto de $\Nat^2$
constituído pelos pares que estão na diagonal, O.s.\@,
\[
  \{\tuple{0,0 }, \tuple{ 1,1 }, \tuple{ 2,2 }, \dots\},
  \]
é a \emph{relação identidade em}~$\Nat$. (Como a relação identidade é
popular, definamos $\Id{A}=\Setabs{\tuple{ x,x }}{x \in
A}$ para qualquer conjunto $A$.) O subconjunto de todos os pares acima da
diagonal, O.s.\@,
\[
  L = \{\tuple{ 0,1 },\tuple{ 0,2 },\ldots,\tuple{ 1,2 },
  \tuple{ 1,3 }, \dots, \tuple{ 2,3 }, \tuple{ 2,4 },\ldots\},
\]
é a relação \emph{menor que}, O.s.\@, $Lnm$ se e só se $n<m$. O subconjunto dos
pares abaixo da diagonal, O.s.\@,
\[
  G=\{ \tuple{ 1,0 },\tuple{ 2,0 },\tuple{
    2,1 }, \tuple{ 3,0 },\tuple{ 3,1 },\tuple{ 3,2 }, \dots\},
\]
é a relação \emph{maior que}, O.s.\@, $Gnm$ se e só se $n>m$. A união
de $L$ com $I$, que podemos chamar $K=L\cup I$, é a relação \emph{menor ou
igual a}: $Knm$ se e só se $n \le m$. Do mesmo modo, $H=G \cup
I$ é a \emph{relação maior ou igual a}. Estas relações
$L$, $G$, $K$, e $H$ são tipos especiais de relações chamadas
\emph{ordens}. $L$ e $G$ gozam da propriedade de que nenhum número está em relação $L$ ou $G$ consigo próprio (O.s.\@, para todo~$n$, não se verifica nem $Lnn$ nem $Gnn$).
Relações com esta propriedade chamam-se \emph{irreflexivas}, e, se
também forem ordens, chamam-se \emph{ordens estritas.}
\end{ex}

\begin{explain}
Embora ordens e identidade sejam relações importantes e naturais, convém
enfatizar que, segundo a nossa definição, \emph{qualquer}
subconjunto de $A^{2}$ é uma relação em~$A$, por mais artificial ou
forçado que pareça. Em particular, $\emptyset$ é uma relação em qualquer
conjunto (a \emph{relação vazia}, que nenhum par de elementos satisfaz), e
$A^{2}$~em si é também uma relação em~$A$ (uma que todo par
de elementos satisfaz), chamada a \emph{relação universal}. Mas também algo como
$E=\Setabs{\tuple{n, m}}{n>5 \text{ ou } m \times n \ge 34}$ conta como
relação.
\end{explain}
  
\begin{prob}
  Liste os !!{element}s da relação $\subseteq$ no conjunto
  $\Pow{\{a, b, c\}}$.
\end{prob}

\end{document}
